% Custom commands, environments and related definitions.
\newcommand\smtlib{SMT-LIB}
\newcommand\cdclt{CDCL(T)}

\newcommand\tool[1]{\textsf{#1}}
\newcommand\tactic[1]{\texttt{#1}}
\newcommand\verit{\tool{veriT}}
\newcommand\cvcfive{\tool{cvc5}}
\newcommand\cvcfour{\tool{CVC4}}
\newcommand\isabelle{\tool{Isabelle/HOL}}

\newcommand\lsymb[1]{\mathbf{#1}}
\DeclareMathOperator*{\subst}{subst}
\DeclareMathOperator*{\reify}{reify}
\DeclareMathOperator*{\bvexplode}{bvexplode}
\newcommand\groundbox[1]{\boxed{#1}}

% Proofs
\newcommand\textAlethe[1]{\texttt{#1}}
\newcommand\ctxsep{$\vartriangleright$}
\newcommand\spctxsep{\multicolumn{1}{|c}{\ctxsep}}
\newcommand\spctx[1]{\multicolumn{1}{|c}{#1}}
\newcommand\spsep{\cline{2-4}}

% Empty lines with dots
\newcommand\aletheLine{ &                       & $\vdots$ & \\}
\newcommand\aletheLineB{& \multicolumn{1}{|c}{} & $\vdots$ & \\}
% with subproof line to the left
\newcommand\aletheLineS {&                       & & $\vdots$ & \\}
\newcommand\aletheLineSB{& \multicolumn{1}{|c}{} & & $\vdots$ & \\}

% Tables for proofs
\newcolumntype{Y}{>{\centering\arraybackslash}X}

\newcommand{\ruleTypeImpl}[1]{%
  \microtypesetup{tracking=false}\textsf{#1}\microtypesetup{tracking=true}%
}
\def\ruleType#1{\ruleTypeImpl{\detokenize{#1}}} % non linked rule (for examples)
\def\proofRule#1{\hyperref[rule:\detokenize{#1}]{\ruleTypeImpl{\detokenize{#1}}}} % linked rule
%\newcommand{\ruleref}[1]{\ruleType{\nameref{rule:#1}}~(\ref{rule:#1})}
\def\ruleref#1{\hyperref[rule:\detokenize{#1}]{\ruleTypeImpl{\detokenize{#1}}}~(\ref{rule:\detokenize{#1}})} % for the rule tables

\NewEnviron{Alethe}{%
\renewcommand\spsep{\cline{2-4}}
\setlength{\arrayrulewidth}{0.8pt}
\addtolength{\tabcolsep}{-4pt}
\begin{xltabular}{\linewidth}{l c Y r}
  \BODY
\end{xltabular}
\addtolength{\tabcolsep}{+4pt}
}
\NewEnviron{AletheS}{%
\renewcommand\spsep{\cline{2-5}}
\setlength{\arrayrulewidth}{0.8pt}
\addtolength{\tabcolsep}{-4pt}
\begin{xltabular}{\linewidth}{l l c Y r}
  \BODY
\end{xltabular}
\addtolength{\tabcolsep}{+4pt}
}

% These avoid the xltabular environment. While xltabular can break on pages
% it seems to interact badly with the RuleDescription environment.
\NewEnviron{AletheX}{%
\renewcommand\spsep{\cline{2-4}}%
\setlength{\arrayrulewidth}{0.8pt}%
\addtolength{\tabcolsep}{-4pt}%
\noindent
\begin{tabularx}{\linewidth}{l c Y r}
  \BODY
\end{tabularx}
\addtolength{\tabcolsep}{+4pt}
}
\NewEnviron{AletheXS}{%
\renewcommand\spsep{\cline{2-5}}%
\setlength{\arrayrulewidth}{0.8pt}%
\addtolength{\tabcolsep}{-4pt}%
\noindent
\begin{tabularx}{\linewidth}{l l c Y r}
  \BODY
\end{tabularx}
\addtolength{\tabcolsep}{+4pt}
}

% Environment for proof-rules
\newcounter{ProofRuleCounter}
\newcommand\currule{\textbf{ERROR}}

\declaretheoremstyle[
headfont=\sffamily\bfseries,%
notefont=\sffamily\bfseries,%
notebraces={}{},%
bodyfont=\normalfont,%
headpunct={},%
headformat={\NAME~\NUMBER:\NOTE},%
break
]{proof-rule-style}

\declaretheorem[name=Rule,style=proof-rule-style,sibling=ProofRuleCounter]{inner-rule}

\makeatletter
\newcommand{\ruleparagraph}{%
  \@startsection{paragraph}{4}%
  {\z@}{1ex \@plus 0.5ex \@minus .2ex}{-1em}%
  {\normalfont\normalsize\bfseries}%
}
\makeatother

\NewEnviron{RuleDescription}[1]{%
\renewcommand\currule{\proofRule{#1}}
\index[rules]{#1}

\begin{inner-rule}[\detokenize{#1}]\label{rule:#1}

\BODY
\end{inner-rule}
}

\theoremstyle{definition}
\newtheorem{RuleExample}{Example}[inner-rule]

% TODO: synchronize colors!
\definecolor{SmtBlue}{HTML}{00007f}
\definecolor{SmtGreen}{HTML}{3b7f31}
\definecolor{SmtStepId}{HTML}{3b7f31}

\newcommand{\grNT}[1]{\textcolor{SmtGreen}{\langle\texttt{#1}\rangle}}
\newcommand{\grT}[1]{\textcolor{SmtBlue}{\texttt{#1}}}
\newcommand{\grRule}{≔}
\newcommand{\grOr}{|}

\newtheorem{example}{Example}
\newtheorem{theorem}{Theorem}[example]
\newtheorem{lemma}{Lemma}[example]
\newtheorem{definition}{Definition}[example]
\newtheorem{seppara}{}

\newpairofpagestyles{titlestyle}{
\ofoot{Version \input{|"./version.sh"}\unskip . Typeset on {\today}.}%
\ifoot{\ccby}%
}
\renewcommand*{\titlepagestyle}{titlestyle}
\setkomafont{publishers}{\small}
